\section{Implicit and Transient Time Integration}
\label{sec:implicit}

\subsection{Steady pseudo-time march}

Steady cases solve $\mathbf{R}(\mathbf{U})=\mathbf{0}$ by a backward-Euler
pseudo-time march. Linearising \cref{eq:fv} about the current iterate gives, at
each pseudo-time step, the linear system
\begin{equation}
  \left(\frac{V_i}{\Delta\tau_i}\mathbf{I}
  +\frac{\partial\mathbf{R}}{\partial\mathbf{U}}\right)\Delta\mathbf{U}
  =\mathbf{R}(\mathbf{U}^m),
  \qquad
  \mathbf{U}^{m+1}=\mathbf{U}^m+\Delta\mathbf{U}.
  \label{eq:steadysystem}
\end{equation}
The pseudo-time step is local to each cell and set from the convective and
viscous spectral radii accumulated during the residual evaluation,
\begin{equation}
  \Delta\tau_i=\frac{\mathrm{CFL}\cdot V_i}
  {\displaystyle\sum_{f\in\partial\Omega_i}\left(|u_n|+a\right)_f\Delta s_f
  \;+\;C_v\sum_{f\in\partial\Omega_i}\frac{\mu_f}{\rho_f}
  \frac{\Delta s_f^2}{V_i}},
  \label{eq:dtau}
\end{equation}
with $C_v=4$ for the viscous contribution. Local time stepping lets the tiny
wall cells and the huge farfield cells advance at their own stable rates, which
is what makes a mesh with a cell-volume range of roughly ten orders of magnitude
tractable at all.

The CFL number is ramped geometrically in log space from
\texttt{cfl\_initial} to \texttt{cfl\_max} over
\texttt{pseudo\_cfl\_ramp\_steps},
\begin{equation}
  \mathrm{CFL}(m)=\mathrm{CFL}_0
  \left(\frac{\mathrm{CFL}_{\max}}{\mathrm{CFL}_0}\right)^{m/M_{\mathrm{ramp}}},
  \label{eq:ramp}
\end{equation}
capped at $\mathrm{CFL}_{\max}$ thereafter. The supplied case values are used
without modification; \cref{tab:controls} lists them.

A residual-based safeguard multiplies the ramped CFL by an adaptive factor. If a
step increases the residual by more than $30\%$, or produces a non-finite
residual, the step is rejected, the previous accepted state is restored, and the
factor is cut to $0.35$ of its value; after a successful step the factor recovers
by $6\%$ up to unity. Rejected-step counts are reported. This is a robustness
device, not an accuracy device: it changes only the path to the steady state.

\subsection{LU-SGS inner solve and the Jacobian approximation}
\label{sec:lusgs}

\Cref{eq:steadysystem} is solved approximately by symmetric Gauss--Seidel sweeps
(LU-SGS); a block-Jacobi variant is also available. The Jacobian is never
assembled. The diagonal is the scalar
\begin{equation}
  D_i=\frac{V_i}{\Delta\tau_i}
  +\frac12\sum_{f\in\partial\Omega_i}\left(|u_n|+a\right)_f\Delta s_f
  +\sum_{f\in\partial\Omega_i}\left(\frac{\mu}{\rho}\right)_f\frac{\Delta s_f^2}{V_i},
  \label{eq:diag}
\end{equation}
where the factor $\tfrac12$ is the upwind split of the face dissipation between
the two adjacent cells. The off-diagonal action on a neighbour is the
scalar-dissipation upwind approximation
\begin{equation}
  \mathbf{O}_{ij}\Delta\mathbf{U}_j=\frac{\Delta s_f}{2}
  \left[\mathbf{A}_j(\mathbf{n})\Delta\mathbf{U}_j
  -\lambda_j\Delta\mathbf{U}_j\right],
  \qquad
  \lambda_j=|u_{n,j}|+a_j+\frac{2\mu}{\rho_j\,\|\Delta\mathbf{x}\|},
  \label{eq:offdiag}
\end{equation}
with the last term present only for viscous runs. Crucially,
$\mathbf{A}_j(\mathbf{n})\Delta\mathbf{U}_j$ is the \emph{exact analytic}
Jacobian-vector product of the Euler normal flux, evaluated matrix-free: no
$4\times4$ block is ever stored. Writing $\mathrm{d}(\rho u_n)=u_n\,\mathrm{d}\rho+\rho\,\mathrm{d}u_n$
and $\mathrm{d}p=(\gamma-1)\left[\mathrm{d}(\rho E)-\tfrac12|\mathbf{u}|^2\mathrm{d}\rho-\rho(u\,\mathrm{d}u+v\,\mathrm{d}v)\right]$,
the product is
\begin{equation}
  \mathbf{A}(\mathbf{n})\,\mathrm{d}\mathbf{U}=
  \begin{bmatrix}
    \mathrm{d}(\rho u_n)\\
    \mathrm{d}(\rho u_n)u+\rho u_n\,\mathrm{d}u+\mathrm{d}p\,n_x\\
    \mathrm{d}(\rho u_n)v+\rho u_n\,\mathrm{d}v+\mathrm{d}p\,n_y\\
    \mathrm{d}u_n(\rho E+p)+u_n\left(\mathrm{d}(\rho E)+\mathrm{d}p\right)
  \end{bmatrix}.
  \label{eq:jacvec}
\end{equation}
This product is verified against central finite differences to $\vJacobianFd$
over $\vJacobianStates$ randomised states (\cref{sec:verification}).

It is worth being precise about what is exact here and what is not, because the
metadata records the solver as \texttt{lu\_sgs\_simplified\_jacobian} and the two
statements must not look contradictory. The \emph{flux Jacobian-vector product}
\cref{eq:jacvec} is exact and analytic. The \emph{overall implicit operator} is
nonetheless a simplified Jacobian: the diagonal \cref{eq:diag} is a scalar rather
than a $4\times4$ block, and the off-diagonal action \cref{eq:offdiag} combines the
exact product with a scalar-dissipation term instead of the true upwind flux
linearisation. That approximation is deliberate and standard for LU-SGS --- it is
what makes the sweeps matrix-free --- and it affects only the rate of convergence to
the steady state, not the steady state itself, because the converged solution is
determined by $\mathbf{R}(\mathbf{U})=\mathbf{0}$ and not by the operator used to get
there.

Each sweep refreshes the ghost increments once, then performs a forward pass over
ascending local cell index and a backward pass over descending index. Between
refreshes the off-rank coupling is frozen, which makes the scheme exactly
block-Jacobi across ranks and pure LU-SGS within a rank --- so the iteration count
can depend mildly on the partition even though the converged answer does not.

The inner iteration is monitored by the true linear residual of the system it is
solving,
\begin{equation}
  r_{\mathrm{lin}}=
  \frac{\left\|\mathbf{R}-\left(D\,\Delta\mathbf{U}+\mathbf{O}\,\Delta\mathbf{U}\right)\right\|_2}
       {\left\|\mathbf{R}\right\|_2},
  \label{eq:linres}
\end{equation}
evaluated with the same matrix-free operator the sweeps use and volume-normalised
like the nonlinear norms. Sweeps are added in small blocks, within the case-file
bounds, until $r_{\mathrm{lin}}$ meets the case target or fails to improve by more
than $2\%$ in three consecutive blocks. The increment is deliberately \emph{not}
reset between blocks --- an early version did reset it, which silently discarded
all previous work and capped the achievable inner convergence.

One negative result is worth recording, because it is a tempting shortcut.
Inflating the diagonal by a factor $\beta>1$ makes the sweeps look dramatically
better: at $\beta=2$ the reported ratio reached $\num{2.2e-16}$. That number is
meaningless, because it measures a different (easier) system than the one being
solved; the true residual ratio of \cref{eq:linres} was $0.82$. The diagonal is
therefore left at its consistent value \cref{eq:diag}, and the slow-but-honest
convergence of SGS at high CFL is accepted instead.

\subsection{Convergence criteria}
\label{sec:criteria}

All reported residual norms are volume-weighted root-mean-square rates of change
of the cell average, globally MPI-reduced:
\begin{equation}
  \|\mathbf{R}\|_2=\left(\frac{\sum_i \sum_k
  \left(R_{i,k}/V_i\right)^2 V_i}{\sum_i V_i}\right)^{1/2},
  \qquad
  \|\mathbf{R}\|_\infty=\max_{i,k}\left|R_{i,k}/V_i\right| .
  \label{eq:norms}
\end{equation}
Dividing by the cell volume converts the integral residual into a rate, which
makes the norm independent of local cell size and therefore comparable across
meshes and rank counts. \Cref{sec:stationarity} describes the termination test
built on these norms, and the defect that had to be fixed in it.

\subsection{Transient dual-time BDF2 for the vortex street}
\label{sec:bdf2}

The $\Reinf=200$ cylinder is integrated in physical time by a genuine two-loop
dual-time scheme. The outer loop steps physical time with second-order backward
differences; the inner loop drives the resulting nonlinear system to convergence
at each physical step.

With $a_0=\tfrac32,\ a_1=-2,\ a_2=\tfrac12$ the BDF2 residual is
\begin{equation}
  \mathbf{R}^\ast(\mathbf{U}^{n+1})=\mathbf{R}(\mathbf{U}^{n+1})
  -\frac{V_i}{\Delta t}\left(a_0\mathbf{U}^{n+1}+a_1\mathbf{U}^{n}+a_2\mathbf{U}^{n-1}\right),
  \label{eq:bdf2}
\end{equation}
and the inner iteration solves $\mathbf{R}^\ast=\mathbf{0}$ with the same LU-SGS
machinery, using the diagonal
\begin{equation}
  D_i^\ast=\frac{V_i}{\Delta\tau_i}+\frac{a_0 V_i}{\Delta t}
  +\frac12\sum_f\left(|u_n|+a\right)_f\Delta s_f+\dots,
  \label{eq:bdf2diag}
\end{equation}
i.e.\ carrying \emph{both} the pseudo-time and the physical-time term, so each
inner iteration is a Newton-like step on the time-accurate system rather than a
pseudo-time march with a source term. The first physical step uses the
self-starting BDF1 coefficients $a_0=1,a_1=-1,a_2=0$.

Two properties required by the benchmark hold explicitly in the implementation:

\begin{enumerate}
  \item $\mathbf{U}^{n}$ and $\mathbf{U}^{n-1}$ are stored in separate arrays and
  are \textbf{frozen throughout all inner iterations} for $\mathbf{U}^{n+1}$.
  Only the $\mathbf{U}^{n+1}$ iterate changes inside the inner loop.
  \item The history is shifted ($\mathbf{U}^{n-1}\!\leftarrow\!\mathbf{U}^{n}$,
  $\mathbf{U}^{n}\!\leftarrow\!\mathbf{U}^{n+1}$) only \textbf{after} the inner
  solve for that physical step has been accepted.
\end{enumerate}

The inner loop terminates when the total transient residual $\|\mathbf{R}^\ast\|_2$
has fallen by the case target relative to its value at the first inner iteration
of that physical step, subject to the minimum iteration count. The convergence
measure is the \emph{total} residual including the physical-time term, as the
benchmark specifies, not the spatial residual alone. Per-step inner counts, target
misses and the final ratio are accumulated and reported in
\texttt{metadata.json}.

The run is declared \texttt{statistically\_periodic} only if it reaches the
requested horizon \emph{and} the inner solve converged on at least $95\%$ of
physical steps. A run that reaches the final time while frequently missing the
inner target is an under-converged dual-time solve and is reported as
\texttt{failed} rather than presented as a usable transient.

\begin{table}[htbp]
  \centering
  \caption{Production run controls. Every value is taken from the supplied case
  file and used unmodified; no case uses a loosened setting.}
  \label{tab:controls}
  \small
  \begin{tabular}{lrrrrr}
    \toprule
    Case group & Max steps & $\mathrm{CFL}_0$ & $\mathrm{CFL}_{\max}$
      & Ramp steps & Inner \\
    \midrule
    NACA inviscid $\Minf=0.15$ & \num{20000} & 1.0 & 100 & \num{2000} & 3--50 \\
    NACA inviscid $\Minf=0.80$ & \num{30000} & 1.0 & 100 & \num{3000} & 3--50 \\
    NACA inviscid $\Minf=2.0$ & \num{40000} & 0.2 & 50 & \num{5000} & 3--80 \\
    NACA laminar $\Minf=0.15$ & \num{40000} & 1.0 & 100 & \num{5000} & 3--80 \\
    NACA laminar $\Minf=0.80$ & \num{40000} & 0.5 & 80 & \num{5000} & 3--80 \\
    NACA laminar $\Minf=2.0$ & \num{50000} & 0.2 & 50 & \num{7000} & 3--100 \\
    Cylinder $\Reinf=20$ & \num{30000} & 1.0 & 100 & \num{3000} & 3--80 \\
    Cylinder $\Reinf=200$ & \num{30000} & 1.0 & 1.0 & 0 & 5--1000 \\
    \bottomrule
  \end{tabular}
\end{table}

For the transient case the supplied production settings are used exactly:
$\Delta t=0.01$, $t_f=300.0$ ($\num{30000}$ physical steps), inner iteration
bounds $5$--$1000$, inner target $10^{-3}$ on the total transient residual, and a
fixed pseudo-time CFL of $1.0$.
